\begin{figure}[htbp]
\centering
\begin{tikzpicture}[
    node distance=0.45cm,
    >=Stealth,
    phaseAnalysis/.style={
        rectangle, rounded corners=4pt,
        draw=black, thick,
        fill=blue!12, align=center,
        minimum width=8cm, minimum height=1.0cm,
        inner sep=4pt
    },
    phaseSpec/.style={
        rectangle, rounded corners=4pt,
        draw=black, thick,
        fill=violet!15, align=center,
        minimum width=8cm, minimum height=1.0cm,
        inner sep=4pt
    },
    phaseBuild/.style={
        rectangle, rounded corners=4pt,
        draw=black, thick,
        fill=green!18, align=center,
        minimum width=8cm, minimum height=1.0cm,
        inner sep=4pt
    },
    phaseEval/.style={
        rectangle, rounded corners=4pt,
        draw=black, thick,
        fill=orange!22, align=center,
        minimum width=8cm, minimum height=1.0cm,
        inner sep=4pt
    },
    phaseSynth/.style={
        rectangle, rounded corners=4pt,
        draw=black, thick,
        fill=gray!18, align=center,
        minimum width=8cm, minimum height=1.0cm,
        inner sep=4pt
    },
    output/.style={
        font=\footnotesize\itshape,
        text=black!65,
        align=left,
        anchor=west
    },
    arrow/.style={->, thick, draw=black!75}
]

% --- Fases ---
\node[phaseAnalysis] (f1) {\textbf{Fase 1:} Mapeamento Sistemático da Literatura (MSL)};
\node[phaseAnalysis, below=of f1] (f2) {\textbf{Fase 2:} Identificação de riscos éticos testáveis};
\node[phaseSpec, below=of f2] (f3) {\textbf{Fase 3:} Especificação das técnicas de \textit{fuzzing}};
\node[phaseBuild, below=of f3] (f4) {\textbf{Fase 4:} Desenvolvimento das \textit{fuzzing engines}};
\node[phaseEval, below=of f4] (f5) {\textbf{Fase 5:} Avaliação empírica em modelos reais};
\node[phaseSynth, below=of f5] (f6) {\textbf{Fase 6:} Documentação, análise e síntese dos resultados};

% --- Saídas como anotações laterais à direita ---
\node[output, right=0.3cm of f1.east] {Saída: 3 princípios prioritários\\(\textit{Fairness}, \textit{Accountability}, \textit{Transparency})};
\node[output, right=0.3cm of f2.east] {Saída: taxonomia de 11 riscos\\(6 implementáveis + 5 não-implementáveis)};
\node[output, right=0.3cm of f3.east] {Saída: mapeamento risco $\rightarrow$ técnica de \textit{fuzzing}};
\node[output, right=0.3cm of f4.east] {Saída: 6 \textit{fuzzing engines} (RF1, RF2, RF4,\\RA2, RT1, RT2) cobrindo os 3 princípios};
\node[output, right=0.3cm of f5.east] {Saída: métricas, falhas detectadas\\e comparação entre provedores};
\node[output, right=0.3cm of f6.east] {Saída: conclusões, recomendações\\e evidências consolidadas};

% --- Setas ---
\draw[arrow] (f1) -- (f2);
\draw[arrow] (f2) -- (f3);
\draw[arrow] (f3) -- (f4);
\draw[arrow] (f4) -- (f5);
\draw[arrow] (f5) -- (f6);

\end{tikzpicture}
\caption{Etapas metodológicas da pesquisa, agrupadas por natureza: análise (azul), especificação (roxo), construção (verde), avaliação (laranja) e síntese (cinza).}
\label{fig:metodologia}
\end{figure}